\paragraph{Theorem E.2.17 (Proved: the full shifted homotopy module).}
Use the right-module convention of E.2.16. Write
\[
 \begin{gathered}
 R=\mathbb ZG,\qquad q_i=h_i-1,\qquad S=\mathbb Z\Gamma,\\
 \Gamma=G/\langle h_1,h_2\rangle
  =\langle a,b,c,d\mid a^3=b^4=c^3=d^4=abcd=1\rangle.
 \end{gathered}
 \tag{E.P1}
\]
where $a,b$ belong to the first complement, $c,d$ to
the second, and $\mu=ab=(cd)^{-1}$. The full integral
second-homotopy module is
\[
 \pi_2(P)\simeq R/q_2R\ \oplus\ R/q_1R
                    \ \oplus\ T_L\oplus T_R.
 \tag{E.P2}
\]
The last two modules are inflated from the following
explicit six-generator presentations over $S$.
Put $N_a=1+a+a^2$, $N_b=1+b+b^2+b^3$, and define
$N_c,N_d$ similarly.

For $T_L$ use generators $p_i,r_i$ for $0\leq i\leq2$,
put $v_i=p_i-r_i$, and define
\[
 \begin{split}
 \Omega_L={}&v_0(1+d-d^2-d^3)
       +v_1(d^3-d)+v_2(d^2+d^3),\\
 p_iN_a={}&0,\qquad r_iN_b=0\quad(0\leq i\leq2),\\
 \Omega_L(a-1)={}&\Omega_L(b-1)
       =\Omega_L(c-1)=\Omega_L(d-1)=0.
 \end{split}
 \tag{E.P3}
\]
For $T_R$ use generators $B_i,C_i$, put $w_i=B_id+C_i$,
and define
\[
 \begin{split}
 \Omega_R={}&w_0a(1-b-b^2+b^3)
              +w_1a(b^2-1)-w_2a(1+b^3),\\
 B_iN_c={}&0,\qquad C_iN_d=0\quad(0\leq i\leq2),\\
 \Omega_R(a-1)={}&\Omega_R(b-1)
       =\Omega_R(c-1)=\Omega_R(d-1)=0.
 \end{split}
 \tag{E.P4}
\]
These are complete presentations: six generators and
ten relations in each block. In particular
\[
 q_1q_2\pi_2(P)=0,\qquad
 \ker(q_1:\pi_2(P)\to\pi_2(P))
  \cap\ker(q_2:\pi_2(P)\to\pi_2(P))=T_L\oplus T_R.
 \tag{E.P5}
\]
This improves the earlier squared annihilator and
specifies its entire central-intersection summand.
Its coinvariants are $(\mathbb Z/12)^6$, so (E.P2)
recovers $\pi_2(P)_G=\mathbb Z^2\oplus(\mathbb Z/12)^6$.

Let $\Pi=\pi_1(Y)$ and let $\Sigma\simeq\pi_1(\Sigma_6)$
be its quotient by $z=h_1h_2$, as in E.2.13. Then
\[
 \pi_2(Y)\simeq(\mathbb Z\Pi)^{24}
          \oplus\operatorname{Res}^{\Gamma}_{\Sigma}(T_L\oplus T_R),
 \qquad
 \pi_2(Y)_\Pi\simeq\mathbb Z^{110}\oplus(\mathbb Z/6)^2.
 \tag{E.P6}
\]
The surface-group summand is inflated to $\Pi$, so
$z$ acts trivially on it. Formulas (E.P2)--(E.P6)
do not identify the smooth type of $P$ or $Y$.

\begin{proof}
\emph{A small chain presentation with the actual peripheral map.}
The calculation in E.2.16 permits the following model
through degree three for the first complement. Its
bases have ranks $(1,2,6,6)$, with degree-one basis
$(a,b)$ and degree-two basis $(r,x,t,y,z,e)$.
The first two differentials are
\[
 d_1=(a-1,b-1),\qquad
 d_2(r)=\begin{pmatrix}N_a\\h_1N_b\end{pmatrix},
 \quad d_2(x)=d_2(t)=d_2(y)=d_2(z)=d_2(e)=0.
 \tag{E.P7}
\]
The six degree-three columns are the six relations
\[
 t(a-1),\quad t(b-1),\quad xq_1,\quad yq_1,\quad zq_1,
                         \quad x(1+b)+3t.
 \tag{E.P8}
\]
The degree-two norm $N_a$ in (E.P7) is written with
the same expression as in (E.P3), but here $a^3=h_1$,
not $1$. The second complement is the identical
model with $(a,b,h_1)$ replaced by $(c,d,h_2)$.

The ordinary one-relator complex $C_K$ is the subcomplex
on degrees zero and one and the cell $r$. It is an
aspherical model: it is the double mapping cylinder
of the degree-three and degree-four maps of circles,
and its universal cover is a tree of contractible
lines and strips. Equivalently it is the standard
Seifert spine of the torus-knot exterior. Projection
of (E.P7)--(E.P8) onto $C_K$ retains the actual
peripheral map. In degree one this map is
\[
 g\longmapsto(N_a,0),\qquad
 u,w,d\longmapsto0,\qquad
 \Theta\longmapsto(b,1).
 \tag{E.P9}
\]
Its degree-two $r$ coefficient is $b-1$ on the swept
cell $g\Theta$ and zero on the other nine cells.
Its remaining components are exactly (E.C19).
These additional columns were tracked in the
E.2.16 certificate, rather than chosen from their
ordinary homology values.

Here and below the small complex is only a
presentation through the degree needed for $H_2$.
Replacing redundant degree-three columns by the six
generators of their image in (E.P8) changes no
degree-two cycle or boundary. Reducing a peripheral
degree-two column modulo that image changes it by
a chain homotopy with values in degree three.
Products with a torus and the mapping cones below
therefore retain the same degree-two homology.
For the comparison with an aspherical target one
may first keep the original, full complex: its
projection to $C_K$ extends by zero in degrees at
least three, since the Fox column in (E.P7) is
injective and hence the $r$ row of $d_3$ is zero.
Only after taking the kernel is its image in degree
two shortened as in (E.P8). Thus no false
acyclicity of a truncated target is used.

\emph{The exact kernel complex for the splice.}
Let $T(s)$ be the four-torus chain complex whose
first circle has deck image $s$ and whose other
three circles have trivial image. Its differential
contracts that first circle and multiplies by $s-1$.
The two vertex complexes for the actual splice are
\[
 V_1=C_{E_1}\otimes T(h_2),\qquad
 V_2=C_{E_2}\otimes T(h_1).
 \tag{E.P10}
\]
We have reversed the two product factors in the
second vertex only for notation; the Koszul signs
below account for that reversal. The boundary
complex $B$ is the nine-torus complex with ordered
coordinates
\[
 (g_1,u_1,w_1,d_1,g_2,u_2,w_2,d_2,\Theta),
 \quad\text{deck images }(h_1,1,1,1,h_2,1,1,1,\mu).
 \tag{E.P11}
\]
Its differential is the usual exterior contraction
by those deck images minus one.

For clarity all source signs can be specified
without an implicit boundary convention. Write
a boundary product cell as $\alpha\wedge\beta$
with an optional final $\Theta$, where $\alpha$
uses the first four coordinates and $\beta$ the
second four. Put $\epsilon=1$ when $\Theta$ occurs
and $0$ otherwise. Its first vertex map applies
the E.2.16 peripheral map to $(\alpha,\Theta)$,
keeps $\beta$, and has sign
$(-1)^{\epsilon(|\alpha|+|\beta|)}$.
The second map applies that peripheral map to
$(\beta,\Theta_2)$, keeps $\alpha$, and has sign
$(-1)^{|\alpha|(|\beta|+\epsilon)+\epsilon|\beta|}$.
If $\epsilon=1$, multiply its coefficient on the
right by $-\mu$. This is the exact cellular map
of $\Theta_1\mapsto\Theta_2^{-1}$, since
$\mu_2=cd=\mu^{-1}$. The suspended cells in (E.C19)
have the interval first; these signs convert from
the ordered exterior cells in (E.P11).
Denote the two maps by $j_1,j_2$.

The chain complex of the splice is the cone of
$(j_1,-j_2):B\to V_1\oplus V_2$.
Project each $C_E$ to $C_K$ and each $T(s)$ to its
first circle, and project $B$ to its three coordinates
$g_1,g_2,\Theta$. The result is the complex of the
aspherical graph of spaces
\[
 (K(K_1,1)\times S^1)\ \cup_{T^3}\
                         (K(K_2,1)\times S^1).
 \tag{E.P12}
\]
It is the model $N$ of E.2.9. Both peripheral maps
on groups are injective by E.2.10; their common
group is $\langle h_1,h_2,\mu\rangle=\mathbb Z^3$.
Consequently its universal cover is a Bass--Serre
tree of contractible vertex and edge spaces and
is contractible. Its chain complex has ranks
$(2,7,9,5,1)$, including its degree-four cell.

The indicated projection is a degreewise split
surjection of the full chain complexes. The
peripheral Fox coefficient in (E.P9) makes the
boundary square commute; higher degrees commute
as well by the product construction and the
zero higher groups of $C_K$. Let $D$ be its kernel.
The target has zero positive-degree homology, so
the long exact sequence gives
$H_2(D)\simeq H_2(\widetilde P)=\pi_2(P)$.
Using the small models for the relevant images,
the ranks of $D$ through degree three are
\[
                         (0,6,40,123).
 \tag{E.P13}
\]
Its six degree-one generators are the unused
regular-factor circles at the two vertices.
The six shifted boundary circles $u_i,w_i,d_i$
map to them by unit entries, because each circle
maps to zero in its own complement and identically
to the opposite regular torus. Cancel these six
unit pairs. This leaves a presentation of $H_2(D)$
with34 generators and123 relations, with no
uncomputed kernel.

Here is also an explicit block definition of that
matrix. In each $V_i$ keep the cells killed by
projection to $C_{K_i}\otimes C(S^1)$; their ranks
in degrees two and three are17 and45. In $B_2$
keep the33 cells containing at least one of
$u_1,w_1,d_1,u_2,w_2,d_2$. In those ordered bases
the matrix is
\[
 \mathcal M=
 \begin{pmatrix}
 (d_{V_1})_3&0&(j_1)_2\\
 0&(d_{V_2})_3&-(j_2)_2
 \end{pmatrix}:
 R^{45}\oplus R^{45}\oplus R^{33}
             \longrightarrow R^{17}\oplus R^{17}.
 \tag{E.P14}
\]
Every entry is specified by (E.P7)--(E.P11) and
(E.C19); (E.P14) itself is a complete finite
presentation of the desired module. The formula
includes the actual derived peripheral comparison.
In particular a splitting of the two cusp terms
has not been assumed from their localized homology.

\emph{Reduction to the compact presentation.}
The exact certificate
\path{code/maps_splice_module.py} checks both
peripheral chain equations and $d_2d_3=0$, cancels
the indicated unit pairs, and reduces (E.P14) to
16 generators and52 nonzero columns. It computes
in the actual amalgam
\[
 G=(K_1\times\mathbb Z)\mathop{*}_{\mathbb Z^3}
                                      (\mathbb Z\times K_2),
 \tag{E.P15}
\]
not in a finite quotient. To specify its normal
form, the torus-knot factor uses $h^r$ times an
alternating word in $a,a^2$ and $b,b^2,b^3$.
Modulo the peripheral subgroup, strip initial
$ab$ and $b^3a^2=(ab)^{-1}$, allowing their last
letter to be a partial power; choose $b^3$ in
the sole length-one tie with $a$. Inverting these
left coset representatives gives right coset
representatives. The usual alternating amalgam
normal form, followed by a final element of
$\mathbb Z^3$, is unique. These elementary word
operations are implemented exactly, including
both central exponents.

In the resulting16-generator presentation two
generators occur only in their respective single
relations $q_2e_1=0$ and $q_1e_2=0$. They give the
first two direct summands of (E.P2). For the other
14 generators, two generators $t_0,t_1$ have
trivial action of $a,b,c,d$. There are three
generators initially killed by $q_2$ whose pure
relations include $h_1N_b$; the identity
$h_1N_b(b-1)=1-h_1$ kills them by $q_1$ as well.
There are three initially killed by $q_1$ with
pure $N_c$ relations, and three more with pure
$h_2N_d$ relations. They are killed by $q_2$.
The remaining generators already have both
central relations. Thus the whole14-generator
summand is a module over $S$, without imposing
an additional quotient. The pure norm columns
are exact differences of the recorded columns
and are checked before setting either central
element equal to1.

For a readable final reduction, the14 generators
separate into two seven-generator blocks. The
first has generators $U_0,U_1,U_2,X,Y,Z,t_1$,
with $t_1$ invariant, $U_iN_b=0$, and
\[
 \begin{split}
 (U_0+X-Yd+Zd)N_a&=0,\\
 (U_1-Y(1+d)+Z)N_a&=0,\\
 (U_2+X+2t_1+Z(d-1))N_a&=0,\\
 X(1+d)+3t_1&=0.
 \end{split}
 \tag{E.P16}
\]
The second has generators $B_i,C_i,t_0$, with
$B_iN_c=C_iN_d=0$, $t_0$ invariant, and the single
additional relation $t_0=\Omega_R$. These formulas
follow by collecting the displayed16-generator
columns; the certificate retains every column
and its unit reductions.

In the first block put
\[
 v_0=X-Yd+Zd,\quad v_1=-Y(1+d)+Z,
 \quad v_2=X+2t_1+Z(d-1),\qquad
 r_i=U_i,\quad p_i=U_i+v_i.
 \tag{E.P17}
\]
The linear part of this change on $(X,Y,Z)$ is
\[
 \begin{pmatrix}1&0&1\\-d&-1-d&0\\d&1&d-1\end{pmatrix},
 \quad\det=1,
\]
over $\mathbb Z[d,d^{-1}]$. It is therefore an
integral group-ring basis change. Its inverse gives
\[
 X=v_0(1-d^2)+v_1(d^2-d)+(v_2-2t_1)d^2.
\]
Substitute in the last relation of (E.P16), using
the invariance of $t_1$. The result is
$t_1=\Omega_L$. Eliminate $t_1$ with this unit
coefficient. Its four invariance relations become
the four relations in (E.P3). Eliminating $t_0$
in the second block gives (E.P4). This proves the
compact presentation in (E.P2)--(E.P4), including
the integral central-intersection information.

For the two cyclic axis summands, the other
central element is a non-zero-divisor: their
group rings are free over its Laurent polynomial
ring. Thus neither contributes to the common
kernel of $q_1,q_2$. This proves (E.P5). On
coinvariants, the four invariance relations in
each block vanish, while the norm relations
become three copies of multiplication by3 and
three copies of multiplication by4. Each block
therefore has coinvariants $(\mathbb Z/12)^3$.
This is an independent recovery of E.2.11.

\emph{Restriction to the actual circle fibre.}
The universal covers of $P$ and $Y\times\mathbb R$
are $\Pi$-equivariantly homotopy equivalent by
E.2.13. The map $\Pi\to G/\langle h_i\rangle$
is injective, since $\chi(h_i)=\pm12$, and has
index12. Hence each of the first two group-ring
summands in (E.P2) is free of rank12 over
$\mathbb Z\Pi$. The last two factor through the
index12 surface subgroup $\Sigma<\Gamma$ and
have trivial $z$ action. This proves the module
formula in (E.P6).

For its exact coinvariants use the regular
12-sheet permutation representation with
exponents $(4,-3,-4,3)$ on $(a,b,c,d)$.
Evaluating the compact12-by20 presentation gives
an integral144-by240 matrix. Its Smith form has
56 unit entries and two entries6; its free
cokernel rank is86. The certificate computes
the same result from the independently retained
14-by26 presentation. Adding the24 free
coinvariants proves the last formula in (E.P6).
All these are integral module and homotopy
statements. No smooth surgery or smooth
classification is inferred from them.
\end{proof}
